% Generated draft: requires mathematical and source review.
\documentclass[11pt]{article}
\usepackage[margin=25mm]{geometry}
\usepackage{fontspec}
\setmainfont{Latin Modern Roman}
\usepackage{amsmath,amssymb,mathtools,unicode-math}
\setmathfont{Latin Modern Math}
\usepackage{enumitem,needspace,xurl,hyperref}
\setlength{\parindent}{0pt}
\setlength{\parskip}{5pt}
\setlength{\emergencystretch}{3em}
\newcommand{\N}{\mathbb{N}}
\newcommand{\Z}{\mathbb{Z}}
\newcommand{\Q}{\mathbb{Q}}
\newcommand{\R}{\mathbb{R}}
\newcommand{\C}{\mathbb{C}}
\newcommand{\F}{\mathbb{F}}
\newcommand{\A}{\mathbb{A}}
\newcommand{\PP}{\mathbb{P}}
\newcommand{\E}{\mathbb{E}}
\newcommand{\eps}{\varepsilon}
\newcommand{\dd}{\,\mathrm{d}}
\begin{document}
\section*{UnsolvedMath ID: 30004537}
\textbf{Source problem code:} catalogue entry
\subsection{Definitions and mathematical statement}
Let \(f\) denote the Gaussian stationary process mentioned in the source, and let \(P_f(N)\in[0,1]\) denote its persistence probability at horizon \(N\). The source does not define the persistence event, the parameter set or value domain of \(f\), or the admissible horizons \(N\). Consequently, \(P_f(N)\) remains a source-dependent quantity rather than a fully specified probability in this draft.

\textbf{Requested statement.} Determine conditions on \(f\) that ensure the existence of the persistence exponent
\[
\theta_f=-\lim_{N\to\infty}\frac{1}{N}\log P_f(N),
\]
where the limit must be taken over the intended admissible horizons. The logarithm is the natural logarithm.

For the interpretation requiring a finite exponent and assuming \(P_f(N)>0\) for all sufficiently large admissible \(N\), the requested conclusion can be expressed precisely as
\[
\exists\,\theta_f\in[0,\infty)\quad
\forall\,\varepsilon>0\quad
\exists\,N_0>0\quad
\forall\text{ admissible }N\ge N_0,\qquad
\left|-\frac{\log P_f(N)}{N}-\theta_f\right|<\varepsilon.
\]
Equivalently, under those qualifications,
\[
P_f(N)=\exp\bigl(-\theta_f N+o(N)\bigr).
\]
These are conditional formulations: the source does not specify eventual positivity, whether \(\theta_f=0\) is permitted, or whether an infinite exponent is allowed. If zero persistence probabilities are permitted, an extended-real interpretation would also require an explicit convention for \(-\log 0\).

\textbf{Unresolved scope.} The persistence event must be supplied, including its threshold, strict or non-strict inequality, and observation times. Discrete observations and persistence throughout a continuous interval define different problems; the latter also requires an appropriate measurability framework. Centering, variance normalization, degeneracy exclusions, and spectral or covariance hypotheses are not stated. The source asks for conditions ensuring existence but does not specify whether it seeks sufficient conditions or a necessary-and-sufficient characterization. Its displayed limit describes an exponential logarithmic rate; it does not by itself request a multiplicative asymptotic equivalent or a prefactor.
\subsection{Short English statement}
Determine assumptions under which the persistence probability of a stationary Gaussian process has a well-defined exponential decay rate. The source leaves the persistence event and process framework unspecified, so the statement remains conditional on those definitions. The supplied dated curation labels the problem partially solved; this draft does not independently verify that assessment.
\subsection{Sources}
\begin{itemize}
\item Supplied local UnsolvedMath registry record, permanent ID \(30004537\), title \emph{Existence of Gaussian Persistence Exponents}, collection \emph{Oberwolfach Reports Open Problems}, source problem code \texttt{catalogue entry}. The record lists proposed year \(2020\) and identifies no individual proposer.
\item Supplied report citation: \emph{Stochastic Processes under Constraints} (2021), Oberwolfach Reports, DOI \texttt{10.4171/owr/2020/32}, pp.~1601--1656. \url{https://doi.org/10.4171/owr/2020/32}. The background attributes the problem to workshop open-problem discussions collected in Oberwolfach Reports; no narrower page or contribution is supplied.
\item Supplied literature reference: \emph{Gaussian persistence exponents}. \url{https://arxiv.org/abs/2112.04820}.
\item Supplied literature reference: \emph{Exact log-asymptotics}. \url{https://arxiv.org/abs/2605.20587}.
\item Literature assessment checked at \texttt{2026-08-22}: this is supplied dated curation, not independently verified. It labels the problem \texttt{partially\_solved} and states that existence of the persistence exponent limit is proved under broad spectral-density and positivity hypotheses, while necessary and sufficient conditions and the general exact-asymptotic problem remain unresolved. The supplied background calls this a dated literature triage derived from OWR curation and requests independent verification of time-sensitive status and cited claims.
\item The registry labels the statement \texttt{corrected\_verified}, but its accompanying verification note says only formula delimiters, notation, grammar, and typography were repaired while mathematical content and scope were preserved. This is not an independent mathematical verification of the formulation or literature assessment.
\end{itemize}
\end{document}
